% Setup - do not change
\documentclass[11pt]{article}
\usepackage[top=0.9in, left=0.9in, bottom=0.9in, right=0.9in]{geometry} 
\usepackage{parskip}

\usepackage[english]{babel}
\usepackage[utf8]{inputenc}
\usepackage{amsmath,amsthm,amssymb,graphicx,pdfpages,lipsum,hyperref}
\usepackage[none]{hyphenat}
\usepackage{csquotes}
\usepackage{booktabs,tabularx}
\usepackage{placeins}

\setlength\parindent{0pt}
%%%%%%%%%%%%%%%%%%%%%%%%%%%%%%%%%%%%%%%%%%%%%%%%%%%%%%%%%%%%%%%%%%%
% add other packages here if required
\hypersetup{hidelinks}

%% Bibliography are specified in this file. You can also choose inline bib style if you want to. But make sure your citation style is consistent (and proper)
% For more details on citation: https://library.unimelb.edu.au/recite
\usepackage[sorting = none]{biblatex}
\addbibresource{references.bib}

%%%%%%%%%%%%%%%%%%%%%%%%%%%%%%%%%%%%%%%%%%%%%%%%%%%%%%%%%%%%%%%%%%% the '%' symbol denotes comments

% Begin document creation
% DELETE THE \lipsum PLACEHOLDERS WHEN YOU BEGIN
\title{\textbf{Forecasting Borough-Level Late-Arrival Share for Completed
WAV-Requested HVFHV Trips in New York City}}
\author{
Jiale Xu \\
Student ID: 1600983 \\
%% Replace the link with your github repo
% 1. Remember to escape underscore (\_) in the link.
% 2. Remember to include the commit you want to submit in the link
\href{https://github.com/MAST30034-AppliedDataScience/project-1-individual-LavenderLe7/tree/984b3928c6646017ebeeac3b41c1e6479c681964}{Github repo with commit}
}

\begin{document}
\maketitle

\section{Introduction}
Timely arrival matters to passengers who rely on wheelchair-accessible vehicle (WAV) service.
Late arrivals are also a service-quality concern for platforms whose licensed bases are
required under New York City Taxi and Limousine Commission (TLC) rules to provide equivalent
service to passengers requesting WAVs \cite{nyctlc_fhv_accessibility_base}. At the start of
each four-hour period, only information then available is used to forecast late share for
each pickup borough over the next four hours. The study defines a trip as late when the
vehicle arrives more than 10 minutes after the requested pickup time; late share is the
proportion of completed WAV-requested high-volume for-hire vehicle (HVFHV) trips classified
as late. The forecasts are intended to help operations and
dispatch teams at high-volume for-hire service (HVFHS) platforms prioritise borough-periods
for human review. The forecasts apply only to completed WAV-requested
trips and do not estimate regulatory compliance or profitability.

The study uses TLC HVFHV Trip Record Data from April to December 2025, spanning spring to
early winter, with later months for chronological evaluation. These records provide
the WAV-request and timing information needed to construct the target. Federal holidays
provide calendar context known in advance; Central Park Local Climatological Data
(LCD) observations describe realised weather retrospectively, while archived
High-Resolution Rapid Refresh (HRRR) forecasts provide formal issue-time weather predictors.
The analysis first characterises temporal and spatial late-share
patterns, then evaluates whether a grouped binomial generalised linear model (GLM) and
weighted XGBoost forecast late share more accurately than a matured 24-hour history
benchmark. Both models use the same issue-time information set, comprising location and time,
calendar context, mature service history and HRRR inputs. The design applies a four-hour outcome-maturity
buffer and conservatively assumes that HRRR runs become available three hours after initialisation.

% You can have \section{}, \subsection{}, and \subsubsection{}, \section*{}, \subsection*{}, and \subsubsection*{}
\section{Datasets}
Nine monthly TLC HVFHV Trip Record Data files \cite{nyctlc_trip_records}
covering April--December 2025 are the primary trip-level source. The WAV request flag
identifies WAV-requested trips; \texttt{request\_datetime} records the requested pickup time;
\texttt{on\_scene\_datetime} records arrival at the pickup location;
\texttt{pickup\_datetime} records trip pickup; and \texttt{PULocationID}
supplies the pickup Taxi Zone \cite{nyctlc_hvfhv_dictionary}. A non-missing
pickup timestamp was used to establish completion, and these five fields
support target construction and borough assignment.

To complement the TLC records, the full-year NOAA Local Climatological Data
(LCD) file for the NY City Central Park station (USW00094728) provides realised
temperature, wind, visibility and precipitation for retrospective descriptive
analysis \cite{noaa_lcd}. Formal weather inputs instead come from archived NOAA
HRRR forecasts accessed through dynamical.org \cite{dynamical_hrrr}. NOAA/NCEP
documents the CONUS HRRR at 3-km resolution, with 48-hour forecasts for the
00/06/12/18 UTC cycles \cite{noaa_hrrr}.

The U.S. Office of Personnel Management's 2025 Federal Holidays
\cite{opm_holidays} provide known-in-advance calendar context; eight of 11
NYC-applicable dates fall within the study interval. The TLC Taxi Zone lookup
and shapefile \cite{nyctlc_taxi_zones} support borough assignment, mapping and
selection of the representative HRRR points, but are not separate analytical
datasets. Table~\ref{tab:data_sources} lists the input shapes of the main
analytical sources; the HRRR entry is the standardised local subset rather than
the full remote archive.

\begin{table}[htbp]
    \centering
    \small
    \renewcommand{\arraystretch}{1.08}
    \begin{tabularx}{\textwidth}{@{}
        >{\raggedright\arraybackslash}p{0.52\textwidth}
        >{\centering\arraybackslash}p{0.25\textwidth}
        >{\centering\arraybackslash}X@{}}
        \toprule
        \textbf{Dataset} &
        \textbf{Instances} &
        \textbf{No. Fields} \\
        \midrule
        TLC HVFHV Trip Record Data
            & 183,307,678 trip records
            & 25 \\
        NOAA Local Climatological Data (LCD)
            & 11,401 station reports
            & 125 \\
        Archived NOAA HRRR borough-point subset
            & 270,970 forecast rows
            & 15 \\
        2025 Federal Holiday Calendar
            & 11 holiday dates
            & 8 \\
        \bottomrule
    \end{tabularx}
    \caption{Input shapes of the TLC and external datasets used in the study.}
    \label{tab:data_sources}
\end{table}
\FloatBarrier

\section{Preprocessing}
\textbf{TLC trips and target construction.}
Using the five TLC fields described above, the nine monthly files were filtered
cumulatively as follows:

\begin{itemize}
    \setlength\itemsep{0pt}
    \setlength\parskip{0pt}
    \setlength\parsep{0pt}
    \item \textbf{WAV scope:} 543,349 of the 183,307,678 input records reported
    that the passenger requested a WAV.
    \item \textbf{Timestamps:} 543,349 records were retained; none lacked a
    request, arrival or pickup timestamp.
    \item \textbf{Study interval:} 543,347 requested pickups fell within April--
    December 2025 (2 excluded).
    \item \textbf{Pickup geography:} 543,305 records mapped through the TLC
    Taxi Zone lookup to a valid pickup zone in one of the five boroughs; 42
    records that did not map were excluded \cite{nyctlc_taxi_zones}.
\end{itemize}

For trips with non-missing \texttt{pickup\_datetime}, define
\begingroup
\setlength{\abovedisplayskip}{4pt}
\setlength{\belowdisplayskip}{4pt}
\begin{equation*}
    \begin{gathered}
        \mathrm{arrival\_offset}_i
            = \texttt{on\_scene\_datetime}_i
            - \texttt{request\_datetime}_i
            \quad \text{(minutes)}, \\
        \texttt{late\_10}_i
            = \mathbf{1}\{\mathrm{arrival\_offset}_i > 10\}.
    \end{gathered}
\end{equation*}
\endgroup
The TLC Data Dictionary defines \texttt{request\_datetime} as the time the
passenger requested to be picked up; therefore, it was treated as requested
pickup time, not app-submission time \cite{nyctlc_hvfhv_dictionary}.
Negative offsets represented early arrivals, while large positive offsets
remained observed outcomes; neither group was trimmed.

For pickup borough \(b\) and four-hour period \(t\), the forecasting target was
\begingroup
\setlength{\abovedisplayskip}{2pt}
\setlength{\belowdisplayskip}{2pt}
\begin{equation*}
    \mathrm{late\_share}_{b,t}
        = \frac{\mathrm{late\_trip\_count}_{b,t}}
        {\mathrm{wav\_request\_count}_{b,t}},
        \qquad \mathrm{wav\_request\_count}_{b,t} > 0.
\end{equation*}
\endgroup

\textbf{External data.}
The high-level processing steps were:

\begin{itemize}
    \setlength\itemsep{0pt}
    \setlength\parskip{0pt}
    \setlength\parsep{0pt}
    \item \textbf{Observed LCD:} Eight fields were loaded. Of the 11,401 input
    reports, 11,024 routine FM-15 and FM-16 surface observations were retained,
    while daily and monthly summaries were excluded
    \cite{noaa_surface_report_types}. NOAA records LCD times in local standard
    time \cite{noaa_lcd_documentation}; conversion to New York civil time left
    8,260 study-period reports. Hourly observations were consolidated using median
    temperature, maximum wind, minimum visibility and event indicators. These hourly
    values were then aggregated into 1,650 four-hour summaries using the corresponding
    mean, extreme and event measures. Missing hours were not imputed, and 1,606 periods
    had complete hourly coverage.
    \item \textbf{Archived HRRR:} The nearest forecast grid point to each
    projected borough centroid was fixed. From the standardised 270,970-row
    archive subset \cite{dynamical_hrrr}, issue-time-eligible valid times inside
    each target period produced one complete six-feature forecast row for all
    8,250 borough-periods; no imputation was required.
    \item \textbf{Calendar and geography:} The Taxi Zone lookup
    \cite{nyctlc_taxi_zones} assigned borough labels and supplied polygons for
    spatial analysis and HRRR point selection. OPM holidays within the study
    interval \cite{opm_holidays} marked all corresponding four-hour periods.
\end{itemize}

\textbf{Panel construction.}
WAV outcomes, citywide request counts and descriptive LCD summaries were
left-joined to the complete 8,250-row grid of 1,650 periods by five boroughs.
Zero-request periods retained an undefined late share, not zero lateness. The
April--October target covered 6,172 of 6,420 rows (96.14\%). No source join
removed a grid row: the base panel contained 8,250 rows and 36 fields, and
the six HRRR summaries were joined one-to-one to form the 42-field modelling
panel.

Forecast-time predictors excluded contemporaneous requests and outcomes, as
well as realised LCD weather; the modelling section specifies historical
maturity and HRRR issue-time eligibility.

\section{Analysis and Geospatial Visualisation}
Exploratory analysis was restricted to April through October, with November and
December reserved for chronological evaluation. The development sample comprised
428,633 completed WAV-requested trips across the complete 6,420-row
borough-period grid, with a pooled late share of 9.42\%.

\subsection{Distribution and temporal structure}
Arrival offsets were asymmetric: the median was 4.30 minutes and the 95th
percentile was 12.18 minutes, while 19.35\% were negative. Only 24 trips
exceeded 60 minutes, although the maximum was 212.4 minutes; these rare extremes
were retained as observed outcomes. Zero-request cells retained an undefined
late share rather than being imputed.

Monthly completed-request volume remained within 57,857--64,300, whereas late
share peaked at 13.64\% in June and fell to 5.57\% in October. This divergence
between broadly similar support and changing risk motivates chronological
evaluation and explicit temporal context.

Figure~\ref{fig:borough_time_risk} shows that time-of-day patterns were
borough-specific rather than a common citywide cycle. The Bronx late share
peaked at 16.2\% from 04:00--07:59, whereas Queens peaked at 21.3\% from
20:00--23:59 on 10,062 requests. Staten Island's rates ranged from 12.4\%
to 18.8\% on only 112--2,753 requests per time block;
the borough also contained all 248 zero-request periods. These patterns
motivate representing borough--time context jointly and accounting for request
support in model fitting and evaluation, since rates based on small denominators
may be unstable.

\begin{figure}[htbp]
    \centering
    \includegraphics[width=0.97\textwidth]{../../outputs/figures/wav_borough_time_risk.pdf}
    \caption{Late-arrival share and completed WAV-request volume by borough and time block.}
    \label{fig:borough_time_risk}
\end{figure}

\FloatBarrier

\subsection{Geospatial and observed-weather evidence}
Figure~\ref{fig:zone_risk} reveals Taxi Zone-level late-share heterogeneity
that borough aggregation would otherwise smooth over, while separating spatial
risk from request support. The late-share panel displays only 241 of the 262
zones meeting the 100-request rule; together, these zones represented 99.85\% of development
requests. JFK Airport had both the largest support (8,355 requests) and the
highest late share among displayed zones (44.15\%). Therefore, it is a
high-support, high-risk
hotspot rather than an unstable rate driven by a very small denominator. The
remaining zones are masked to avoid ranking poorly supported rates. This map
is descriptive; the formal forecasting unit remains the
borough by four-hour period.

\begin{figure}[!htbp]
    \centering
    \includegraphics[width=0.97\textwidth]{../../outputs/figures/wav_taxi_zone_support_and_risk.pdf}
    \caption{Development-period completed WAV-request support and late share by
    Taxi Zone.}
    \label{fig:zone_risk}
\end{figure}

\FloatBarrier

Among borough-periods with observed targets and complete LCD hourly coverage,
the study-defined adverse-weather flag (any precipitation, minimum visibility
below 5 km, or maximum wind speed at least 10 m/s) had an 11.64\% late share
across 80,424 requests, compared with 8.96\% across 336,180 requests in other
periods. The 2.68 percentage-point difference is an unadjusted association based
on realised Central Park observations used as a citywide proxy. It neither
identifies a causal weather effect nor establishes that forecast weather improves
prediction. Because realised target-period LCD observations would not be
available at forecast issue, LCD variables are excluded from the model
predictors.

The observed borough-time heterogeneity and month-to-month variation motivate a
chronological forecasting design. Therefore, the models test whether calendar and location
context and mature service history predict late share, and whether issue-time
HRRR adds value beyond them under chronological evaluation.

\FloatBarrier

\section{Modelling}
\subsection{Forecast Design and Feature Engineering}
\begin{sloppypar}
For each borough \(b\) and target period \(t\), both predictive models estimated
\(\widehat p_{b,t}\), the forecast late share during the four hours beginning at issue
time. Historical risk over an eligible set of source periods \(W\) was pooled as
\begin{equation*}
    r_b(W)=\frac{\sum_{s\in W}L_{b,s}}{\sum_{s\in W}N_{b,s}},
\end{equation*}
where \(L_{b,s}\) and \(N_{b,s}\) denote late\hyp trip and completed\hyp request counts.
Historical aggregates used completed\hyp trip periods ending at least four hours earlier,
so the most recent eligible period was lag 2. Four such features captured different
temporal scales: lag-2 risk, matured 24-hour and seven-day histories, and an expanding
same-slot history for the same borough, weekday and time block. Count pooling prevented
sparse source periods from receiving equal weight. A development-period audit found a
maximum request-to-completion delay of 212.4 minutes, below the 240-minute maturity
buffer.
\end{sloppypar}

Calendar predictors captured weekday, holiday, time block, linear trend and smooth
annual seasonality. Log-transformed lag-2 borough demand, matured 24-hour borough
demand and lag-2 citywide demand represented recent request volume.

The latest HRRR run was eligible only if \(\text{initialisation time}+3\text{ hours}\leq\text{issue time}\).
The latest eligible run was fixed for the target period and summarised into mean temperature, mean
relative humidity, total precipitation, any rain, mean wind speed and maximum wind
gust. Current-period information, realised LCD observations and later HRRR runs were
excluded. Indicators distinguished unavailable historical risk features from observed
zero risk.

Models were fitted on April--August, validated on September--October, then refitted on
April--October and evaluated with fixed parameters on November--December. At every
issue time, predictors excluded the target period, and earlier outcomes entered
subsequent histories only after satisfying the maturity rule. Preprocessing parameters
were estimated from fitting data only. The primary metric was request-weighted MAE in
percentage points; secondary diagnostics comprised RMSE, calibration assessments and
expanding monthly backtests.

\subsection{Matured 24-Hour History Benchmark}
The matured 24-hour pooled rate---the count-weighted late share across source
periods \(t-7,\ldots,t-2\)---was the benchmark, with training-period borough
and then global fallbacks. It was a transparent comparator and potential
fallback; GLM and XGBoost were the predictive models.

\subsection{Grouped Binomial GLM}
As a denominator-aware, interpretable specification, a binomial GLM modelled
grouped late and on-time counts \cite{mast30027_binomial_regression}:
\begingroup
\setlength{\abovedisplayskip}{2pt}
\setlength{\belowdisplayskip}{2pt}
\begin{equation*}
\begin{aligned}
    L_{b,t}\mid N_{b,t} &\sim
        \operatorname{Binomial}(N_{b,t},p_{b,t}),\\
    \operatorname{logit}(p_{b,t})
        &=\beta_0
          +\beta^{(\text{borough}\times\text{time block})}_{b,h(t)}\\
        &\quad+\beta^{(\mathrm{weekday})}_{w(t)}
          +\beta_{\mathrm{holiday}}\,
          \mathbf{1}\!\left\{\operatorname{date}(t)\in\mathcal H\right\}
          +\boldsymbol{x}_{b,t}^{\mathsf T}\boldsymbol{\beta}_x.
\end{aligned}
\end{equation*}
\endgroup
Here \(p_{b,t}\) is the conditional expected late share; \(h(t)\), \(w(t)\) and
\(\mathcal H\) denote time block, weekday and holiday dates; and
\(\boldsymbol{x}_{b,t}\) contains the retained standardised non-categorical
predictors. Reference coding represented joint borough--time-block and weekday
effects. Target-period \(N_{b,t}\) entered the grouped response, not the forecast
inputs. The model assumed a common conditional late probability within each
borough-period and additive log-odds effects; coefficients were conditional
associations, not causal effects.

\subsection{Weighted XGBoost}
As a flexible contrast, XGBoost used the same issue-time information set but
represented borough and time block separately. It fitted \(L_{b,t}/N_{b,t}\),
with the contemporaneous \(N_{b,t}\) weighting the loss rather than entering as
a forecast input. Its logistic objective bounded forecasts to the unit interval,
while regularised depth-two trees captured nonlinearities and interactions
\cite{chen_guestrin_xgboost}; chronological validation selected the tree count
before refitting. Unlike the GLM, it imposed no additive mean structure. Both
models relied on sufficiently stable predictor--outcome relationships and on the
continued validity of the maturity and HRRR availability assumptions.

\section{Results and Discussion}
\subsection{Predictive Value Beyond Mature History}
In the November--December evaluation, both learned models provided modest
average predictive value beyond mature history, but gains were uneven.
Table~\ref{tab:model_comparison} shows that weighted XGBoost had the lowest MAE,
2.589 percentage points---12.4\% below the benchmark. Evaluation RMSE preserved
the ordering, so this lead was not specific to MAE. Because the GLM and XGBoost
shared the issue-time information set, XGBoost's further gain is consistent with
its more flexible nonlinear and interaction structure rather than access to
different feature families.

\begin{table}[htbp]
    \centering
    \small
    \renewcommand{\arraystretch}{1.08}
    \begin{tabularx}{\textwidth}{@{}
        >{\raggedright\arraybackslash}p{0.27\textwidth}
        >{\centering\arraybackslash}p{0.13\textwidth}
        >{\centering\arraybackslash}p{0.13\textwidth}
        >{\centering\arraybackslash}p{0.13\textwidth}
        >{\centering\arraybackslash}X@{}}
        \toprule
        \textbf{Method} &
        \textbf{Validation MAE (pp)} &
        \textbf{Evaluation MAE (pp)} &
        \textbf{Evaluation RMSE (pp)} &
        \textbf{Evaluation calibration gap (pp)} \\
        \midrule
        Matured 24-hour history & 2.968 & 2.955 & 4.130 & -0.010 \\
        Grouped Binomial GLM & 2.607 & 2.776 & 3.928 & +0.445 \\
        Weighted XGBoost & \textbf{2.550} & \textbf{2.589} &
            \textbf{3.716} & +0.045 \\
        \bottomrule
    \end{tabularx}
    \caption{Request-weighted predictive performance on the validation and evaluation
    sets. Evaluation comprised 1,731 scored borough-periods and 114,672 completed 
    WAV-requested trips. Calibration gap is predicted minus observed request-weighted
    late share.}
    \label{tab:model_comparison}
\end{table}

Although adverse observed weather was associated with higher late share, the
HRRR ablation showed that incremental forecast value varied between validation
and evaluation: refitting XGBoost without the six HRRR summaries lowered
validation MAE from 2.550 to 2.508 but raised evaluation MAE from 2.589 to 2.664.
In validation, permuting the mature same-slot, 24-hour and seven-day late-risk features
increased MAE by 0.577, 0.272 and 0.074 pp, respectively. The largest positive HRRR importance,
for forecast mean temperature, was only 0.015 pp. XGBoost was therefore
most sensitive to mature service history. These permutation results reflect predictive
reliance rather than causal effects, and correlated predictors may prevent individual
feature contributions from being cleanly separated.


\subsection{Reliability and Operational Constraints}
\label{sec:reliability}
Aggregate calibration did not imply borough-period accuracy: the benchmark's
gap was closer to zero than XGBoost's despite its higher MAE, while the GLM had
the largest positive gap. Across four one-month-ahead development backtests,
XGBoost exceeded the benchmark MAE by 0.30 pp in July and 0.99 pp in August,
but was about 0.40 pp lower in both September and October. Its advantage was
therefore not temporally uniform.

Errors were largest where completed-request support was sparse. XGBoost's
evaluation MAE was 10.195 pp for Staten Island (1,748 requests), versus
2.471 pp elsewhere. Table~\ref{tab:support_mae} shows the same pattern across
request-volume bands, with prediction error declining sharply as completed-request
support increased.

\begin{table}[htbp]
    \centering
    \small
    \renewcommand{\arraystretch}{1.08}
    \begin{tabularx}{\textwidth}{@{}
        >{\raggedright\arraybackslash}p{0.39\textwidth}
        >{\centering\arraybackslash}p{0.28\textwidth}
        >{\centering\arraybackslash}X@{}}
        \toprule
        \textbf{Requests per borough-period} &
        \textbf{Evaluation WAV requests} &
        \textbf{Weighted MAE (pp)} \\
        \midrule
        1--10   & 1,851  & 9.803 \\
        11--50  & 15,513 & 4.402 \\
        51--100 & 23,415 & 2.770 \\
        $>$100  & 73,893 & 1.970 \\
        \bottomrule
    \end{tabularx}
    \caption{XGBoost evaluation error by completed-request support.}
    \label{tab:support_mae}
\end{table}

Therefore, request weighting makes the headline MAE more representative of high-volume
periods but may mask substantially larger errors in sparse-support groups.

Two retrospective screens illustrate the concentration--coverage trade-off. First,
a fixed 10\% predicted late-share screen covered 6.09\% of requests and captured
11.46\% of late trips. Second, a volume-aware ranking multiplied predicted late
share by the expanding same-slot mean completed-request count; its top 10\% contained
31.09\% of requests and 35.22\% of late trips. The risk-only screen concentrated
risk more sharply but reached fewer late trips, whereas volume-aware ranking reached
more late trips with only modest enrichment. Because the rules selected different
shares of periods, the results show a trade-off rather than a capacity-matched policy
comparison.

\FloatBarrier

\section{Recommendations}
When review and response capacity is limited, operations and dispatch teams at HVFHS platforms should use
forecasts, with analytics support, to prioritise borough-periods for review. Live platform data should
then determine whether and how the platform responds. Alerts for sparse-support
borough-periods should receive additional scrutiny and should not by themselves trigger an
operational response. Therefore, we recommend two operational responses and a controlled pilot.

\textbf{Offer targeted driver incentives to reinforce WAV coverage.}
When the forecast prioritises a borough-period, the HVFHS platform's dispatch team should review
outstanding WAV requests, available WAV-capable vehicles, cancellations and known disruptions.
If live data show that available coverage cannot meet outstanding requests, the team should
offer drivers eligible to serve WAV requests location- and time-specific incentives to accept
WAV trips in that borough or reposition there during the four-hour period. The team should
monitor coverage, end the offer when the shortfall clears and check that the intervention
is not creating new service pressure elsewhere. If live data do not confirm a coverage problem,
the team should record the alert without offering incentives or changing resources. This
focuses operational attention and incentives on prioritised borough-periods, not platform-wide.

\textbf{Prioritise recovery for outstanding WAV requests.}
For each prioritised borough-period, the dispatch team should review outstanding WAV requests
ahead of the usual order and check their live assignment status. The forecast prioritises
review; it does not identify which individual request will be late or allocate a vehicle.

If live data show that a request remains unassigned, loses an assigned vehicle or is affected
by a known disruption, the team should confirm that it remains active, then re-offer it to
drivers eligible to serve WAV requests or rematch it with an available WAV-capable vehicle.
If ordinary matching cannot restore the trip, the team should follow a pre-approved backup
WAV-service pathway and track the recovery until its outcome is recorded. Requests progressing
normally should remain under usual dispatch. The forecast prioritises the borough-period, but
live request status triggers recovery; a borough-wide coverage shortfall is not required.

\textbf{Run a platform-specific pilot before wider rollout.}
The platform's operations and analytics teams should rebuild the late-share forecast with the
platform's own historical data and test both responses in a small pilot. This staged rollout is
warranted because the current model pooled completed trips across platforms, achieved only a
modest average gain over the history benchmark, varied in relative performance over time and
produced much larger errors in sparse borough-periods, especially Staten Island.

Before the pilot, the platform should define review and response capacity, alert routing,
activation and stopping conditions, and team responsibilities. During the pilot, it should
record alerts, actions and uptake; staff time and operating costs; service outcomes, including
completed-trip lateness, cancellations and unserved requests; and spillovers to other boroughs
or periods.

After outcomes mature, the operations and analytics teams should compare XGBoost with the
matured 24-hour history benchmark, compare the predicted late-share and volume-aware rankings
under equal review capacity, and compare forecast-informed responses with usual manual review
under equal review and response capacity. The platform should expand the responses only if
prospective results show better WAV service outcomes at acceptable additional cost without
material deterioration in sparse-support groups or elsewhere; otherwise, it should restrict,
revise or stop them.

\section{Conclusion}
This report investigated whether borough-level late share for completed WAV-requested HVFHV
trips could be forecast over the following four hours using only information available at issue
time. Both learned models outperformed the matured 24-hour history benchmark in the November–December
evaluation, with weighted XGBoost performing best, although its advantage was not temporally uniform
and prediction errors were substantially larger in sparse-support borough-periods. These findings
support evaluating forecasts as a tool to prioritise operational review rather than automate vehicle
allocation. Finally, we recommend using forecast-prioritised borough-periods to guide reviews for targeted
driver incentives and earlier recovery of outstanding WAV requests. Live platform data should determine whether
either response is required. These responses should first be evaluated in a platform-specific controlled pilot before
wider rollout.

\clearpage

% BEGIN REFERENCES SECTION
\printbibliography

\end{document}
